\documentclass[12pt]{article}
\makeatletter
\def\input@path{{../../style/}}
\makeatother

\usepackage{../../style/psetconfig}
\title{Homework 3}
\author{Songyu Ye}
\date{\today}

\begin{document}
\psettitle

\begin{problem}[1 (Hatcher AT, 4.1.3)]
For an H-space $(X,*)$ with multiplication
\[
\mu \colon X\times X \to X,
\]
show that the natural group operation on $\pi_n(X,*)$ can also be defined by the rule
\[
(f+g)(x)=\mu(f(x),g(x)).
\]

Recall that an $H$-space is a space $X$ together with a point $e\in X$ and a continuous map $\mu\colon X\times X\to X$ such that $\mu(e,x)$ and $\mu(x,e)$ are homotopic to $x$ through $(X,e)$-based homotopies, meaning that the homotopies fix the basepoint $e$ for all $x\in X$. The point $e$ is called the unit of the H-space, and $\mu$ is called the multiplication. An immediate consequence of the definition is $\mu(e,e) = e$.
\end{problem}

\begin{solution}
Represent classes in $\pi_n(X,e)$ by maps $f,g:I^n\to X$ with $f|_{\partial I^n}=g|_{\partial I^n}=e$.

Define the usual sum by concatenation in the last coordinate:
\[
(f+g)(u,t)=
\begin{cases}
f(u,2t), & 0\le t\le \tfrac12,\\
g(u,2t-1), & \tfrac12\le t\le 1,
\end{cases}
\]
and define a second operation by
\[
(f\star g)(u,t)=\mu(f(u,t),g(u,t)).
\]
Since $\mu(e,e)=e$, this is again a based map.

These two operations satisfy the interchange law
\[
(f_1\star g_1)+(f_2\star g_2)=(f_1+f_2)\star(g_1+g_2),
\]
which follows immediately from the definitions.
Moreover, the constant map at $e$ is a unit for $\star$ up to based homotopy, since
\[
\mu(f,e)\simeq f,\qquad \mu(e,f)\simeq f
\]
through based maps, and it is also the unit for $+$.
Hence the Eckmann--Hilton argument implies that $+$ and $\star$ agree.
Therefore
\[
[f]+[g]=[\mu\circ(f,g)].
\]
\end{solution}

\begin{problem}[2 (Hatcher, 4.1.20)]
Show that $[X,Y]$ is finite if $X$ is a finite CW complex and $\pi_i(Y)$ is finite for $i\le \dim X$.
\end{problem}

\begin{solution}
We argue by induction on the skeleta of $X$. Let $X^k$ denote the $k$-skeleton.
Since $X$ is a finite CW complex, only finitely many skeleta occur, and $X=X^m$ for $m=\dim X$.

For $k=0$, $X^0$ is a finite set of points, so
\[
[X^0,Y]\cong \pi_0(Y)^{\# X^0},
\]
which is finite because $Y$ has finitely many components: indeed $\pi_0(Y)$ is finite since $\pi_i(Y)$ is finite for $i\le m$, in particular for $i=0$.

Assume $[X^{k-1},Y]$ is finite. The space $X^k$ is obtained from $X^{k-1}$ by attaching finitely many $k$-cells. For a fixed map
\[
f\colon X^{k-1}\to Y,
\]
its extensions over one $k$-cell correspond to nullhomotopies of the attaching map in $Y$, and whenever extensions exist, the set of extension classes is a torsor under $\pi_k(Y)$. Since only finitely many $k$-cells are attached, the set of possible extensions of $f$ to $X^k$ is finite because $\pi_k(Y)$ is finite.

Thus each class in $[X^{k-1},Y]$ has only finitely many lifts to $[X^k,Y]$. Since the base set $[X^{k-1},Y]$ is finite by induction, it follows that $[X^k,Y]$ is finite. Inducting up to $k=m$ gives that $[X,Y]$ is finite.
\end{solution}

\begin{problem}[3]
Determine $\pi_1(\R\P^n)$ and $\pi_n(\R\P^n)$. Use this to show that:
\begin{enumerate}[(a)]
    \item there is no retraction of the inclusion $\R\P^k \subset \R\P^n$;
    \item the action of $\pi_1(\R\P^n)$ on $\pi_n(\R\P^n)$ is trivial if $n$ is odd and nontrivial when $n$ is even.
\end{enumerate}
\end{problem}
\begin{solution}
For $n\ge 2$, the antipodal map on $S^n$ defines a two-sheeted covering
\[
p:S^n\to \R\P^n.
\]
Hence
\[
\pi_1(\R\P^n)\cong \Z/2\Z.
\]

More generally, if $q:\widetilde X\to X$ is a covering map, then for every $j\ge 2$ the induced map
\[
q_*:\pi_j(\widetilde X)\xrightarrow{\ \cong\ }\pi_j(X)
\]
is an isomorphism. Indeed, every based map $S^j\to X$ lifts uniquely to $\widetilde X$ once a basepoint upstairs is chosen, since $S^j$ is simply connected for $j\ge 2$; the same holds for homotopies because $S^j\times I$ is also simply connected. Applying this to $p:S^n\to \R\P^n$, we get
\[
\pi_j(\R\P^n)\cong \pi_j(S^n)\qquad (j\ge 2).
\]
In particular,
\[
\pi_n(\R\P^n)\cong \pi_n(S^n)\cong \Z.
\]

We now prove the two assertions.

\begin{enumerate}[(a)]
    \item Suppose $k<n$ and there were a retraction
    \[
    r:\R\P^n\to \R\P^k
    \]
    of the inclusion
    \[
    i:\R\P^k\hookrightarrow \R\P^n.
    \]
    Then
    \[
    r\circ i=\id_{\R\P^k},
    \]
    so functoriality gives
    \[
    r_*\circ i_*=\id_{\pi_k(\R\P^k)}.
    \]
    Thus $i_*$ must be injective.

    But by the covering-space fact above,
    \[
    \pi_k(\R\P^k)\cong \pi_k(S^k)\cong \Z,
    \]
    while
    \[
    \pi_k(\R\P^n)\cong \pi_k(S^n)=0
    \]
    since $k<n$ and $S^n$ is $(n-1)$-connected. Therefore
    \[
    i_*:\pi_k(\R\P^k)\to \pi_k(\R\P^n)
    \]
    is a map
    \[
    \Z\to 0,
    \]
    which is not injective. This is a contradiction. Hence no such retraction exists.

    \item The action of $\pi_1(\R\P^n)$ on $\pi_n(\R\P^n)$ is induced by deck transformations of the universal cover
    \[
    p:S^n\to \R\P^n.
    \]
    Since
    \[
    \pi_1(\R\P^n)\cong \Z/2\Z,
    \]
    the nontrivial element acts by the nontrivial deck transformation, namely the antipodal map
    \[
    a:S^n\to S^n,\qquad a(x)=-x.
    \]
    Under the identification
    \[
    \pi_n(\R\P^n)\cong \pi_n(S^n)\cong \Z,
    \]
    this action is therefore multiplication by the degree of $a$. We know that $\deg(a)=(-1)^{n+1}$ because $-I$ acting on $S^n \subset \R^{n+1}$ has determinant $(-1)^{n+1}$. Hence the action is multiplication by $+1$ if $n$ is odd and by $-1$ if $n$ is even. Hence the action is trivial when $n$ is odd and nontrivial when $n$ is even.
\end{enumerate}
\end{solution}

\begin{problem}[4]
Let $* \in A \subset X$ and let $P^*_A X$ be the space of paths in $X$ starting at the basepoint $*$ and ending in $A$. Note the inclusion
\[
j\colon \Omega X \hookrightarrow P^*_A X,
\]
and the endpoint projection
\[
p\colon P^*_A X \to A.
\]
Show from the definitions that:
\begin{enumerate}[(a)]
    \item $\pi_{n-1}(P^*_A X)\cong \pi_n(X,A)$ for $n>0$;
    \item the induced map
    \[
    j_* \colon \pi_{n-1}(\Omega X)\to \pi_{n-1}(P^*_A X)
    \]
    agrees with the map
    \[
    \pi_n(X)\to \pi_n(X,A);
    \]
    \item the induced map
    \[
    p_* \colon \pi_{n-1}(P^*_A X)\to \pi_{n-1}(A)
    \]
    agrees with the connecting map
    \[
    \partial
    \]
    in the long exact homotopy sequence for $(X,A)$.
\end{enumerate}

\medskip

\noindent\textit{Comment.}
The sequence of maps
\[
\Omega X \xrightarrow{\,j\,} P^*_A X \xrightarrow{\,p\,} A
\]
is an example of a fibration. We will see that fibrations also lead to long exact homotopy sequences. In this case, we have a principal fibration with structure group $\Omega X$: in a homotopical sense, $A$ is the quotient of $P^*_A X$ by $\Omega X$. Morally, quotienting out by $\Omega X$ forgets all about the path, except for its endpoint in $A$.
\end{problem}

\begin{solution}
We use the adjunction between maps into a path space and maps out of a cylinder.
A based map
\[
S^{n-1}\to P_A^*X
\]
is the same thing as a map
\[
F\colon S^{n-1}\times I\to X
\]
such that $F(s,0)=*$ for all $s$ and $F(s,1)\in A$ for all $s$. Collapsing $(S^{n-1}\times \{0\})\cup (\{s_0\}\times I)$
to the basepoint identifies this with a based map of pairs
$(D^n,S^{n-1})\to (X,A)$.
Passing to homotopy classes gives
\[
\pi_{n-1}(P_A^*X)\cong \pi_n(X,A),
\qquad n>0.
\]
This proves (a).

For (b), under this identification, a based loop in $\Omega X$ is a map
$S^{n-1}\to \Omega X$
which adjoints to a based map $D^n/S^{n-1}=S^n\to X$
Composing with
\[
j\colon \Omega X\hookrightarrow P_A^*X
\]
just regards such a loop as a path whose endpoint is the basepoint, viewed as a point of $A$. Therefore the induced map
\[
j_*\colon \pi_{n-1}(\Omega X)\to \pi_{n-1}(P_A^*X)
\]
corresponds exactly to the natural map
\[
\pi_n(X)\to \pi_n(X,A).
\]

For (c), under the same adjunction, the endpoint projection
\[
p\colon P_A^*X\to A
\]
sends a path to its endpoint. Thus on a representative $F\in \pi_n(X,A)$, which is a map
\[
F\colon (D^n,S^{n-1})\to (X,A)
\]
of a relative class, the induced map records the restriction
\[
F|_{S^{n-1}}\colon S^{n-1}\to A.
\]
But this is exactly the connecting morphism
\[
\partial\colon \pi_n(X,A)\to \pi_{n-1}(A)
\]
in the long exact sequence of the pair. Hence $p_*$ agrees with $\partial$.
\end{solution}

\begin{problem}[5]
Recall that an element $\alpha\in \pi_n(X,*)$ is a based homotopy class of based maps $S^n\to X$.
\begin{enumerate}[(a)]
    \item Choose a representative map $f$ for $\alpha$, and check that the induced map
    \[
    f_* \colon H_n(S^n;\Z)\to H_n(X;\Z)
    \]
    does not depend on the choice of $f$.
    \item Check that the Hurewicz map
    \[
    h\colon \pi_n(X,x)\to H_n(X),
    \]
    which sends $\alpha\mapsto f_*(1)$, is a group homomorphism.
    \item Show that if $\alpha,\beta\in \pi_n(X)$ are related by the action of $\pi_1$, then $h(\alpha)=h(\beta)$. Reformulate this in terms of free homotopy classes of maps $S^n\to X$.
\end{enumerate}
\end{problem}

\begin{solution}
\begin{enumerate}[(a)]
    \item If $f,g\colon S^n\to X$ are based homotopic, then the induced maps on singular homology are equal. Therefore $f_*=g_*$ on
    \[
    H_n(S^n;\Z),
    \]
    so $f_*(1)$ depends only on the based homotopy class $\alpha=[f]$.

    \item We show that
    \[
    h\colon \pi_n(X,*)\to H_n(X;\Z),
    \qquad
    h([f])=f_*(1),
    \]
    is additive. The group law on $\pi_n(X,*)$ is defined using the pinch map
    \[
    q\colon S^n\to S^n\vee S^n.
    \]
    If $\alpha=[f]$ and $\beta=[g]$, then
    \[
    \alpha+\beta=[(f\vee g)\circ q].
    \]
    On homology, $q_*$ sends the fundamental class $[S^n]$ to the sum of the fundamental classes of the two wedge summands. This can be easily verified by observing that $H^n(S^n\vee S^n;\Z) \cong H^n(S^n;\Z) \oplus H^n(S^n;\Z)$ and writing $q_*(1) = (a,b)$, and then reasoning about the the degree of the projections $S^n\vee S^n \to S^n$ to conclude that $a=b=1$. Hence
    \[
    h(\alpha+\beta)
    =((f\vee g)\circ q)_*(1)
    =(f\vee g)_*\bigl(q_*(1)\bigr)
    =f_*(1)+g_*(1)
    =h(\alpha)+h(\beta).
    \]

    \item The action of $\pi_1(X,*)$ on $\pi_n(X,*)$ can be understood in the cubical model as follows. Represent
    $\alpha\in\pi_n(X,x_0)$
    by
    $F:(I^n,\partial I^n)\to (X,x_0)$.
    Write a point as $(u,t)\in I^{n-1}\times I$. For each fixed $u$,
    $t\mapsto F(u,t)$
    is a loop at $x_0$.

    Let $\gamma$ be a loop at $x_0$. The $\pi_1$-action is
    $(\gamma\cdot F)(u,-)=\gamma * F(u,-) * \gamma^{-1}$. Then $F$ and $\gamma\cdot F$ are freely homotopic via the map which slurps up the conjugation by $\gamma$. It does not fix the boundary, so it is a free homotopy, not a based one. Singular homology is invariant under free homotopy, so if $\alpha$ and $\beta$ lie in the same $\pi_1$-orbit, then
    \[
    h(\alpha)=h(\beta).
    \] This is because a homotopy between $\alpha$ and $\beta$ induces a chain homotopy between the induced maps on singular chains, hence an equality on homology.
    
    This result is equivalent to saying that for $n\ge 2$ the Hurewicz map factors through free homotopy classes of maps $S^n\to X$:
    \[
    \pi_n(X,*)/\pi_1(X,*)\longrightarrow H_n(X;\Z).
    \]
\end{enumerate}
\end{solution}
\end{document}